@@M0@@
gives, by adjunction, a homomorphism:
@@M1@@
and we check immediately on each stalk that it is an isomorphism.

\sgaStableTarget{sga6.e06.para.2-3-2}2.3.2. Now let @@M2@@ be a quasi-coherent graded @@M3@@-Algebra, @@M4@@ a graded @@M5@@-Module, and @@M6@@. For every affine open @@M7@@ of @@M8@@, and every section @@M9@@ of @@M10@@ over @@M11@@, homogeneous of degree @@M12@@, we write @@M13@@ for the homogeneous component of degree @@M14@@ of the Module of fractions @@M15@@; we therefore have:
@@M16@@
If @@M17@@ is an affine open of @@M18@@, then @@M19@@ is the restriction to @@M20@@ of @@M21@@.

We consider the basis of open sets of @@M22@@ formed by the affine opens:
@@M23@@
We write @@M24@@ for the structural morphism @@M25@@, and @@M26@@ for the affine morphism:
@@M27@@
We then put:
@@M28@@
with the notation of 2.3.1. We shall show that the @@M29@@ glue as @@M30@@ and @@M31@@ vary, thereby defining an @@M32@@-Module denoted @@M33@@. The gluing for variable @@M34@@ and fixed @@M35@@ being evident, we are reduced (cf. EGA II 2.5.2) to showing that, if @@M36@@ is a section of @@M37@@ over @@M38@@, homogeneous of degree @@M39@@, then the restriction of @@M40@@ to @@M41@@ is canonically isomorphic to @@M42@@. We consider:
